\documentclass[11pt,a4paper]{article}
\usepackage[utf8]{inputenc}
\usepackage[T1]{fontenc}
\usepackage[a4paper,top=23mm,bottom=21mm,left=23mm,right=23mm]{geometry}
\usepackage{booktabs,tabularx}
\usepackage{xcolor}
\usepackage{fancyhdr}
\usepackage{microtype}
\usepackage{hyperref}

\definecolor{azul}{HTML}{1D4560}
\definecolor{verde}{HTML}{176A63}
\definecolor{cinza}{HTML}{4C5D68}
\hypersetup{
  colorlinks=true,linkcolor=azul,urlcolor=verde,pdflang={pt-BR},
  pdftitle={Enchentes no Rio Grande do Sul - Relatório do dashboard},
  pdfauthor={Gustavo Fidelis, Leonardo Fagundes, Leonardo Mallet e Samuel Bitdinger}
}
\setlength{\parindent}{0pt}
\setlength{\parskip}{6pt}
\setlength{\headheight}{15pt}
\pagestyle{fancy}
\fancyhf{}
\fancyhead[L]{\small\color{cinza}PUCRS · Coleta, Preparação e Análise de Dados}
\fancyhead[R]{\small\color{cinza}INCT SiMAI}
\fancyfoot[L]{\small\color{cinza}Relatório do dashboard}
\fancyfoot[R]{\small\thepage}
\renewcommand{\headrulewidth}{0.3pt}
\renewcommand{\arraystretch}{1.2}

\begin{document}

{\color{verde}\small\bfseries PROJETO EXTENSIONISTA · FASE 3}\par
{\LARGE\bfseries\color{azul} Enchentes no Rio Grande do Sul}\par
{\large Relatório do dashboard}\par
{\small Gustavo Fidelis, Leonardo Fagundes, Leonardo Mallet e Samuel Bitdinger\par
Prof.\textsuperscript{a} Katherine Bianchini Esper de Vargas · Novembro de 2026}

\section*{\color{azul}Introdução}
O \href{http://vhlab.com.br/INCT/floods/}{Flood do INCT SiMAI} oferece uma simulação de
prédios atingidos por bairro de Porto Alegre conforme sobe o Guaíba. O enunciado do
projeto aponta lentidão e instabilidade no acesso à ferramenta. O painel concentra
o nível selecionado, os totais da cidade e o detalhamento por bairro em uma tela. A
segunda tela acrescenta informações sobre danos, interrupções de serviços e
deslocamentos na enchente de 2024.

\section*{\color{azul}Tela 1 · Simulação do Guaíba}
\textbf{Objetivo.} Mostrar quais bairros de Porto Alegre seriam atingidos e estimar
quantos prédios e pessoas seriam afetados em cada nível simulado.

\textbf{Como usar.} Selecione um nível entre 3,0 e 10,0\,m, em passos de 0,1\,m.
Os três indicadores mostram pessoas estimadas, prédios afetados e bairros com pelo
menos um prédio afetado. O ranking mostra até dez bairros e pode ser ordenado por
população estimada ou percentual de prédios afetados.

\textbf{Escolha das visualizações.} As barras horizontais mostram todos os bairros
atingidos no cenário selecionado, com espaço para os nomes e acesso à contagem de
prédios no detalhe de cada barra. A curva usa os 71 níveis da simulação e destaca
o ponto selecionado. O recorde de maio de 2024, 5,35\,m, fica marcado na curva. A
estimativa nesse ponto usa interpolação linear entre 5,3 e 5,4\,m.

\begin{center}
\small
\begin{tabular}{@{}lrrr@{}}
\toprule
\textbf{Nível} & \textbf{Pessoas estimadas} & \textbf{Prédios afetados} & \textbf{Bairros atingidos}\\
\midrule
3,0\,m & 32.119 & 14.652 & 24\\
5,3\,m & 136.389 & 53.989 & 36\\
10,0\,m & 271.886 & 101.219 & 54\\
\bottomrule
\end{tabular}
\end{center}
{\small\color{cinza} Cenários de referência calculados a partir do CSV tratado do INCT SiMAI.}

\textbf{Limitações.} A estimativa de pessoas multiplica a população do bairro
pela proporção de prédios afetados. Isso supõe moradores distribuídos igualmente
entre todos os prédios, inclusive os não residenciais. Os valores descrevem
cenários simulados. As correções do JSON original estão detalhadas no relatório de
limpeza de dados.

\newpage
\section*{\color{azul}Tela 2 · Impactos da enchente de 2024}
\textbf{Objetivo.} Mostrar o que aconteceu com os domicílios e moradores na área
atingida do RS e explorar como alguns impactos e respostas variaram por faixa de
renda domiciliar.

\begin{tabularx}{\linewidth}{@{}p{0.30\linewidth}X@{}}
\toprule
\textbf{Visualização} & \textbf{Como usar e por que foi escolhida}\\
\midrule
Impactos por faixa de renda &
Selecione um dos seis indicadores. As barras permitem comparar cinco faixas de renda.
A linha pontilhada marca o percentual no total do universo do indicador.\\
\addlinespace
Ocorrências nos domicílios e no entorno &
Alterne entre domicílio e bairro ou ruas próximas. Barras horizontais acomodam os
nomes e permitem comparar a frequência dos relatos. A pergunta admite múltiplas respostas.\\
\addlinespace
Onde os moradores passaram a noite &
Mostra o percentual entre os domicílios com alguém que passou ao menos uma noite fora.
A cor mais clara sinaliza estimativas com coeficiente de variação acima de 15\%.\\
\addlinespace
Maio de 2024 nas duas fontes &
O cartão do INCT mostra a parcela estimada da população de Porto Alegre em prédios
atingidos a 5,35\,m. Ao lado, as barras da PEERS mostram ocorrências relatadas
por domicílios na área pesquisada do RS. O evento de 2024 conecta as duas telas.\\
\bottomrule
\end{tabularx}

\textbf{Limitações.} A PEERS é uma pesquisa amostral por telefone. As tabelas
utilizadas apresentam estimativas agregadas da área pesquisada. O painel mantém
separadas as unidades das duas fontes, pessoas em Porto Alegre na simulação e
domicílios da área pesquisada do RS na PEERS. As barras de erro mostram
intervalos aproximados de 95\%. Nas taxas por renda, a aproximação usa o
coeficiente de variação do numerador e não incorpora a incerteza do denominador.

\section*{\color{azul}Execução e fontes}
Após instalar as dependências de \texttt{requirements.txt}, execute
\texttt{python -m streamlit run app/dashboard.py} na raiz do repositório. Use as
abas no topo para alternar entre as telas. O painel carrega os CSVs tratados
presentes no repositório, o que dispensa novas requisições às fontes durante o uso.

\textbf{INCT SiMAI.} \href{http://vhlab.com.br/INCT/floods/}{Flood e simulação de
inundação por bairro}. JSON fornecido à equipe. Prédios do OpenStreetMap e
população do Censo 2022/IBGE.

\textbf{IBGE.}
\href{https://ftp.ibge.gov.br/Meio_Ambiente/PEERS/2024/Abrangencia_xlsx.zip}{Pesquisa
Especial sobre as Enchentes de 2024 no Rio Grande do Sul (PEERS)}. Tabelas 1, 4,
5, 6, 10, 11, 17 e 19, coletadas em 09/11/2026. A data consta nos metadados do
download registrados no repositório.

\section*{\color{azul}Uso de inteligência artificial}
O código das duas telas do dashboard, implementadas em Python com Streamlit,
foi gerado pelo ChatGPT (OpenAI) a partir dos requisitos do enunciado, dos
wireframes e da análise exploratória da equipe.

\end{document}
